% ============================================================================
%  REGRESSION BRACKETING CLOSURE (Wave 4)
%  Validator: theory_v2/val_regression_bracketing_closure.py
% ============================================================================
\section{Regression and general-drift bracketing: closed dichotomy}\label{sec:reg-bracket-closure}

The open row ``fully-general-drift / regression bracketing'' is closed as a
\textbf{dichotomy}, not a universal bracket.

\begin{theorem}[Regression bracketing dichotomy]\label{thm:reg-bracket-dichotomy}
\textbf{(Bounded drift; complete iff.)} For squared loss with $\|b\|_\infty\le B$ known,
Proposition~\ref{thm:reg-iff} gives an \emph{exact} knowability boundary:
$\sign\Delta$ is label-free identifiable iff $|U-2T_S|>2BW$, with adversarial tightness at
$b^\star=-B\sign(f_0-f_a)$.

\textbf{(Fully general drift; impossibility.)} Without a drift budget, the paired-benefit sign
is not universally identifiable: Theorem~\ref{thm:imp} and Theorem~\ref{thm:conj1-dichotomy}
apply; matched evidence with $\TV=0$ and opposite $\sign\Delta$ yields minimax error $\tfrac12$.

\textbf{(Noise.)} Proposition~\ref{thm:reg-noise}: the adapt/freeze decision is invariant to
unknown zero-mean label noise; cardinal risk estimation is confounded.

Thus ``fully general drift'' does not admit a universal label-free bracket
(Conjecture~\ref{conj:gen}, resolved negatively); bounded drift admits a complete
\emph{per-family} characterization (Proposition~\ref{thm:reg-iff}).
\end{theorem}
